\section{Discussion}
\label{sec:discussion}

\paragraph{A regional belt, not isolated hotspots} The clearest result is geographic. Merseyside, Lancashire, Greater Manchester, South Yorkshire and the North East produce more players per boy born than the South and East, and London is the one strong producer in the South. The pattern works at the scale of regions and city-regions: Moran's $I$ is high across TTWAs, region effects remove the residual spatial autocorrelation from the models, and the local tests find a gradient of moderately elevated areas rather than single towns that survive a correction for testing every area. A scout looking for one town that produces players would find the belt instead.

\paragraph{Local opportunity, not elite infrastructure} Within regions, boys born closer to any professional club were more likely to reach the top two divisions, and the effect is stronger for those born since 1990. Distance to the elite Category 1 academies made no difference for the cohorts that grew up after their creation. One reading is that the route into professional football starts at the nearest club, through its academy, development centre or scouting network, and that the most promising boys move to elite academies later. A second reading is that scouts see more of the boys who live near them. The data cannot separate the two, and club locations are themselves a product of historic population and industry, so the distance effect is an association.

\paragraph{Deprivation} The apparent link between deprivation and production comes from region: the North is both more deprived and more productive, and within regions the link disappears. This does not show that deprivation is irrelevant for individuals, only that it does not explain differences between areas once region is known.

\paragraph{Comparison with the birthplace effect} Earlier studies found that athletes are over-represented in mid-sized communities and under-represented in the largest cities \cite{cote2006,macdonald2009}. In English football the largest city, London, is over-represented, and excluding the five largest cities does not change the distance effect. The difference may lie in the sport, in the density of English professional clubs, or in the use of births instead of total population as the baseline.

\paragraph{Players born elsewhere} Half of the study players were born outside England, which shows how international the second tier has also become. The fall in Scottish-born players is the largest change between cohorts; it may reflect changes in Scottish football, in English clubs' recruitment, or in work-permit and registration rules, which this study does not examine.
